\documentclass[letterpaper,twocolumn,10pt]{article}

\usepackage[margin=0.75in]{geometry}
\usepackage{graphicx}
\usepackage{amsmath}
\usepackage{amssymb}
\usepackage{booktabs}
\usepackage{multirow}
\usepackage{hyperref}
\usepackage{xcolor}
\usepackage{microtype}
\usepackage{enumitem}
\usepackage{caption}
\usepackage{subcaption}
\usepackage{array}
\usepackage{tabularx}
\hypersetup{
  colorlinks=true,
  linkcolor=blue,
  citecolor=blue,
  urlcolor=blue
}

\graphicspath{{figs/}}

\title{%
  \textbf{CBS-Bootstrapped MAPPO: Scalable Multi-Agent Path Finding\\
  Under Dynamic Obstacles and Natural Language Control}
}

\author{
  Kyaw Linn Khant \\
  Independent Researcher \\
  Singapore \\
  \texttt{klkhant707@gmail.com}
}

\date{}

\begin{document}
\maketitle

\begin{abstract}
Multi-Agent Path Finding (MAPF) is a critical challenge in warehouse automation and
multi-robot systems. Classical solvers such as Conflict-Based Search (CBS) provide
optimal guarantees but are exponential in the number of agents and fundamentally
cannot handle dynamic environments. Fully learned approaches suffer from catastrophic
cold-start: sparse multi-agent rewards leave the policy gradient uninformative for
the first hundreds of thousands of training steps, producing degenerate always-wait
policies. We present \textbf{CBS-Bootstrapped MAPPO}: a three-phase curriculum that
uses CBS as a \emph{decaying reward signal}---not imitation learning---to bootstrap
a transformer-based multi-agent reinforcement learning policy. The CBS bonus anneals
from weight 1.0 to 0.0 over 500k steps; by Phase~C the policy is fully independent,
continues to improve, and generalises beyond the teacher's reach. We evaluate on
procedurally generated warehouse maps across four difficulty levels
(7\texttimes 7 to 20\texttimes 20, 2--12 agents) under lifelong goal assignment.
On static maps, our policy surpasses CBS at all difficulty levels (71\%--95\% vs.\
CBS 0\%--61\%). Crucially, under dynamic obstacles---where CBS plan invalidation
rates reach 50\%--100\%---our policy maintains 79\%--95\% per-agent success rate
while CBS degrades to 0\%--7\%. We further extend the system with a natural language
zone-assignment interface: a local LLM (Qwen2.5-3B via Ollama) parses operator
commands into warehouse zone targets which the pretrained MARL policy executes
collision-free with no retraining. Training runs entirely on commodity Apple Silicon
hardware (~14 hours for 7M steps). Code and pretrained models are available at
\url{https://github.com/KyawLinnKhant/mapf-cbs}.
\end{abstract}

\section{Introduction}

Multi-Agent Path Finding (MAPF) is the problem of routing $N$ agents through a shared
grid environment from their respective starts to goals, subject to the constraint that
no two agents occupy the same cell (vertex conflict) or swap positions (edge conflict)
at any timestep~\cite{sharon2015cbs}. MAPF is the core computational problem
underpinning modern warehouse automation~\cite{ma2017lifelong}, multi-robot
coordination, and air traffic management. Industrial deployments at scale---warehouses
with dozens of robots, hospitals with autonomous delivery fleets, distribution centres
running 24/7---introduce two complications that classical solvers cannot address.

\textbf{Dynamic obstacles.} Human workers, forklifts, and other unmodelled entities
move through the environment at runtime. A plan that was valid when computed may be
blocked within seconds. CBS plans on a static snapshot of the world; when an obstacle
moves into a planned path, the plan is invalidated and must be recomputed from scratch.
Our experiments show CBS plan invalidation rates of 51\%--100\% with just four moving
obstacles, and full CBS failure on expert-level maps (CBS solve time exceeds the
allowed budget even on static maps).

\textbf{Lifelong operation.} In real deployments, agents are continuously assigned new
goals rather than solving a single batch~\cite{ma2017lifelong}. The relevant metric
is sustained \emph{throughput}---goals reached per unit time---not one-shot plan
optimality. Classical solvers optimise for SoC on a single batch; they have no native
mechanism for lifelong operation without constant replanning.

\textbf{Classical MAPF solvers.} CBS~\cite{sharon2015cbs} finds optimal solutions via
a two-level conflict tree search but scales exponentially in the number of conflicts.
MAPF-LNS2~\cite{li2021mapflns2} improves scalability via large neighbourhood search but
retains the static-world assumption. In our experiments, CBS degrades to 0\% success on
expert-level maps (12 agents, 20\texttimes 20 grid) when capped at 10,000 constraint
tree nodes---a generous budget for real-time operation.

\textbf{Learned MAPF approaches.} PRIMAL~\cite{sartoretti2019primal} and
MAPPO~\cite{yu2022mappo} offer $O(1)$-per-agent inference and naturally handle
environmental change, but suffer from catastrophic cold-start. In standard MAPF
reward structures ($+1$ on goal, $-0.5$ on collision, $-0.01$ per step), random
agent initialisation leads to immediate collisions and no positive reward signal.
The policy gradient is dominated by collision penalties with no gradient toward
coordination for the first hundreds of thousands of steps.

\textbf{Our approach.} We bridge these paradigms by using CBS as a
\emph{temporary reward shaper}. During Phase~A (0--100k steps), CBS computes
optimal actions at each step; agents receive a $+0.3$ bonus when their sampled
action matches CBS. During Phase~B (100k--500k steps), this bonus anneals linearly
to zero. By Phase~C (500k+ steps), the policy is entirely RL-driven. This provides
the warm start needed to exit the zero-gradient region while ensuring the policy is
not permanently anchored to CBS solutions---a requirement for Phase~C generalisation
to dynamic environments.

\textbf{Key distinction.} SCRIMP~\cite{scrimp2023} uses CBS-generated trajectories
for imitation learning, maintaining a demonstration buffer from which the policy
clones behaviour. Our CBS bonus is computed at each step from the \emph{current}
policy state, not from a fixed buffer. More importantly, it decays to zero, allowing
the policy to discover coordination strategies (dynamic obstacle avoidance, flow
corridors, adaptive priority passing) that CBS never produces. No demonstration
buffer is maintained.

\textbf{Contributions.} We make three contributions:
\begin{enumerate}[leftmargin=*,itemsep=2pt,topsep=2pt]
  \item \textbf{CBS as annealing reward shaper:} We show that CBS used as a decaying
    action-match bonus---rather than imitation learning---provides stable curriculum
    bootstrapping for MARL on MAPF. The policy surpasses CBS performance at every
    difficulty level within 500k steps and continues improving for 7M total steps.
  \item \textbf{Curriculum trains beyond the teacher:} On environments with dynamic
    obstacles where CBS plan invalidation reaches 51\%--100\%, our policy maintains
    79\%--95\% per-agent success rate, while CBS+dynamic degrades to 0\%--14\%. The
    policy was not explicitly trained for dynamic obstacles; generalisation emerges
    from the local-reactive observation structure combined with learned coordination.
  \item \textbf{LLM zone assignment:} We demonstrate a natural language interface
    where operator commands are parsed by a local LLM (Qwen2.5-3B, fully offline)
    into warehouse zone assignments, executed collision-free by the pretrained policy
    without retraining. The LLM handles only \emph{intent parsing}---not path
    planning---cleanly separating language understanding from execution.
\end{enumerate}

\section{Related Work}

\subsection{Classical MAPF}

CBS~\cite{sharon2015cbs} solves MAPF optimally via a two-level search: a
high-level constraint tree resolves agent conflicts, and a low-level $A^*$ planner
finds single-agent paths subject to constraints. CBS is complete and optimal but
scales exponentially in the number of conflicts---for dense environments with many
agents, the constraint tree becomes intractably large. MAPF-LNS2~\cite{li2021mapflns2}
applies large neighbourhood search to improve scalability at the cost of optimality
guarantees. LaCAM2 and EECBS offer sub-optimal but faster solutions. All classical
solvers assume a static environment and require expensive replanning when dynamic
obstacles invalidate a plan.

\subsection{Learning-Based MAPF}

PRIMAL~\cite{sartoretti2019primal} uses ODrM* demonstrations for imitation followed
by RL fine-tuning, addressing cold-start through expert demonstrations. SCRIMP
(2023)~\cite{scrimp2023} combines CBS-generated demonstrations with PPO in a curriculum
framework, using a fixed demonstration buffer. DHC and SACHA use communication
networks but do not address dynamic obstacles. MAPPO~\cite{yu2022mappo} demonstrated
the effectiveness of PPO with centralised training for cooperative multi-agent games,
but did not apply it to the MAPF cold-start problem or dynamic environments.

Our approach differs from PRIMAL and SCRIMP in a key respect: we use CBS as a
\emph{live reward signal that decays to zero}, not a fixed demonstration dataset.
The policy never clones CBS trajectories; it receives a diminishing bonus for
agreeing with CBS and is free to diverge as the bonus approaches zero. This allows
Phase~C to discover coordination strategies beyond CBS's capability.

\subsection{Reward Shaping for MARL}

Potential-based reward shaping~\cite{ng1999rewardshaping} guarantees policy invariance:
any potential function $F(s, s') = \gamma\Phi(s') - \Phi(s)$ produces the same
optimal policy as the unshaped MDP. Our CBS bonus is not potential-based---it is
action-matched, not state-based---and intentionally decays to zero, breaking the
invariance guarantee by design. This is \emph{required} for Phase~C: a policy that
permanently receives CBS bonuses has no incentive to learn behaviour CBS cannot produce.
The decay schedule ensures that by Phase~C, the RL objective fully governs learning.

\subsection{LLMs for Robotics and Multi-Robot Control}

Recent work uses LLMs as planners for single-robot manipulation
(SayCan, PaLM-E, RT-2) and as high-level task decomposers (SayPlan, LLM-MCTS).
Attempts to use LLMs as direct MAPF solvers (e.g., path generation for multiple
agents) consistently fail for 3+ agents: the combinatorial action space
(5$^N$ joint actions for $N$ agents per step) exceeds LLM in-context reasoning capacity.

We use the LLM exclusively for intent parsing---extracting spatial zone assignments
from natural language---and delegate all path planning and navigation to the
pretrained MARL policy. This clean separation means the LLM's failure modes
(hallucinated zone names, malformed JSON) can be handled by a deterministic
regex fallback without affecting the navigation policy. Ashish Kapoor's
``ChatGPT for Robotics''~\cite{liang2023code4robotics} establishes the principle
that LLMs are better at generating code/intents than direct control; our zone
assignment interface applies this principle to multi-robot coordination.

\subsection{Entropy Collapse in Multi-Agent RL}

Policy entropy collapse---where the policy distribution concentrates on a single
degenerate action---is a known failure mode in policy gradient methods. In
multi-agent settings, the problem is amplified: the optimal response for agent $i$
often degenerates to ``wait for all other agents,'' which minimises collision
penalties but also prevents any goal completion. We observe this failure mode
empirically at step 5.25M (entropy drops to 0.086, goals drop from 27 to 15) and
recover by raising $c_\text{ent}$ from 0.005 to 0.02. We document this as a
reproducible training event with recovery mechanism.

\section{Problem Formulation}

\subsection{MAPF and Lifelong MAPF}

A MAPF instance is a tuple $(G, A, s, t)$ where $G = (V, E)$ is an undirected grid
graph, $A = \{a_1, \ldots, a_N\}$ is a set of $N$ agents, $s: A \to V$ maps agents
to start positions, and $t: A \to V$ maps agents to goal positions. At each
discrete timestep $\tau$, each agent takes an action
$u_i \in \{\textsc{stay},\,\textsc{N},\,\textsc{S},\,\textsc{E},\,\textsc{W}\}$.
A solution is \emph{valid} if no vertex conflicts (two agents occupy the same cell)
or edge conflicts (agents swap positions) occur at any timestep.

\textbf{Cost metrics.} The classical objective minimises \emph{sum-of-costs}
$\text{SoC} = \sum_i |\pi_i|$, where $|\pi_i|$ is the length of agent $i$'s path.
For lifelong MAPF, we use: (1)~\emph{goals reached per episode}---total goal
completions across all agents in an episode of fixed length $T = 512$ steps, and
(2)~\emph{per-agent success rate}---the fraction of agents that reach at least one
goal within the episode.

\textbf{Lifelong operation.} Following~\cite{ma2017lifelong}, each agent is
immediately assigned a new random goal upon reaching its current goal. The environment
thus operates continuously without terminal states (except episode length $T$).
This models real warehouse deployments where robots are never idle.

\subsection{Dynamic Obstacles}

We extend the MAPF environment with $K$ uncontrolled obstacles. At initialisation,
$K$ obstacles are placed on free cells (not overlapping agents, walls, or goals).
Each step, each obstacle moves according to a stochastic pattern sampled uniformly
from: (1)~\emph{random walk}---a uniformly random valid adjacent cell, and
(2)~\emph{patrol}---alternating between two fixed waypoints at distance 4--6.
Obstacles are observable: the agent's $7\times 7$ local observation includes a
dedicated channel with obstacle positions (value 0.5, distinct from other agents
at value 1.0 and walls).

Crucially, dynamic obstacles are \emph{not modelled by CBS}. CBS plans on the static
map snapshot at episode start. Any obstacle that moves into a planned path step
constitutes a \emph{plan invalidation event}. We quantify invalidation rate as the
fraction of planned steps whose target cells are occupied by a dynamic obstacle when
the agent attempts to move there. This metric directly measures CBS's inability to
handle dynamic environments.

\subsection{Observation and Action Space}

Each agent $i$ observes a $7\times 7$ window centred on its current position,
encoded as 7 binary channels:
\begin{enumerate}[leftmargin=*,itemsep=1pt,topsep=1pt]
  \item Walls (static obstacles)
  \item Other agent positions (value 1.0)
  \item Own position (value 1.0, single cell)
  \item Goal direction (4-channel one-hot: N/S/E/W toward goal)
  \item Dynamic obstacle positions (value 0.5)
  \item Goal cell position (value 1.0)
  \item Distance-to-goal encoding (scalar, normalised to $[0,1]$)
\end{enumerate}
This $7\times 7 \times 7 = 343$-channel tensor is flattened to a 150-dimensional
vector (after a spatial pooling layer in the MLP encoder).

The action space is $\mathcal{A} = \{0,1,2,3,4\}$ corresponding to
\{stay, north, south, east, west\}. Actions that would move an agent into a wall or
out of bounds are treated as stay (masked after sampling).

\section{Method}

\subsection{MAPPO with Transformer Communication}

Figure~\ref{fig:arch} shows the full architecture. Each agent $i$ processes its
150D observation through a shared two-layer MLP encoder (hidden dimension 128, ReLU
activations) to produce an embedding $\mathbf{e}_i \in \mathbb{R}^{128}$.

\begin{figure}[t]
  \centering
  \includegraphics[width=\columnwidth]{architecture}
  \caption{MAPPO architecture with transformer communication. The MLP encoder is
    shared across all agents. The transformer uses a padding mask to support
    variable numbers of active agents (padding mask = 0 for inactive agent slots).
    The centralised critic mean-pools all agent embeddings during training only;
    at deployment, only the actor is used.}
  \label{fig:arch}
\end{figure}

\textbf{Transformer communication.} All $N$ embeddings $\{\mathbf{e}_i\}_{i=1}^N$
are processed by a two-layer multi-head self-attention transformer~\cite{vaswani2017attention}
(4 heads, 128-dimensional, pre-norm LayerNorm, dropout 0.0 during evaluation). The
multi-head attention is:
\[
  \text{MHA}(Q, K, V) = \text{softmax}\!\left(\frac{QK^\top}{\sqrt{d_k}} + M\right) V
\]
where $M_{ij} = -\infty$ if agent $j$ is masked (inactive), 0 otherwise.
This padding mask enables variable-$N$ operation: a checkpoint trained with $N=12$
can be evaluated with $N=8$ by masking positions 8--11. This is critical for real
warehouse deployments where some robots may be offline, charging, or newly added.

\textbf{Centralised Training, Decentralised Execution (CTDE).} The context-enriched
embedding $\tilde{\mathbf{e}}_i$ from the transformer feeds two heads:
\begin{itemize}[leftmargin=*,itemsep=2pt,topsep=2pt]
  \item \textbf{Shared actor:} $\pi_\theta(a_i | \tilde{\mathbf{e}}_i)$---a linear
    layer producing logits over 5 actions. Each agent uses its own $\tilde{\mathbf{e}}_i$
    only (decentralised execution).
  \item \textbf{Centralised critic:} $V_\phi(s) = f_\phi(\frac{1}{N}\sum_i \tilde{\mathbf{e}}_i)$---
    a linear layer applied to the mean-pooled embedding over all agents. The critic
    thus sees the full joint state implicitly through the mean-pooled representation.
\end{itemize}

\textbf{Training objective.} Standard PPO clipped surrogate with entropy regularisation:
\begin{equation}
  \mathcal{L}(\theta) = \mathcal{L}_\text{CLIP} - c_\text{ent}\,\mathcal{H}[\pi_\theta]
    + c_v\,\mathcal{L}_V
  \label{eq:ppo}
\end{equation}
with clipping $\epsilon = 0.2$, $c_\text{ent} = 0.02$ (raised from 0.005 at step 5.25M
to recover from entropy collapse), $c_v = 0.5$. Per-agent advantages are computed via
GAE~\cite{schulman2017ppo} ($\lambda = 0.95$, $\gamma = 0.99$) and normalised
across agents per rollout.

\subsection{CBS-Bootstrapped Three-Phase Curriculum}

Figure~\ref{fig:curriculum} illustrates the CBS bonus annealing schedule and difficulty
curriculum progression.

\begin{figure}[t]
  \centering
  \includegraphics[width=\columnwidth]{curriculum}
  \caption{\emph{Left:} CBS bonus $\delta$ annealing over training steps. Phase~A
    maintains $\delta = 0.3$; Phase~B anneals linearly to zero over 400k steps;
    Phase~C has zero CBS dependency. \emph{Right:} Difficulty curriculum progression.
    Rapid advancement (easy$\to$expert by step 164k) is driven by the CBS warm start;
    all transitions occur within Phase~A/B.}
  \label{fig:curriculum}
\end{figure}

\textbf{Phase A (0--100k steps): CBS oracle warm start.}
At each step $\tau$, the CBS oracle computes the optimal joint action $a^*_i$ for all
agents given the current grid state (walls, agent positions, goals). If agent $i$'s
sampled action $a_i$ matches $a^*_i$, it receives an additional shaped reward:
\[
  r^\text{CBS}_i = \delta_0 \cdot \mathbf{1}[a_i = a^*_i], \quad \delta_0 = 0.3
\]
The CBS computation uses a node-capped variant
($\text{max\_ct\_nodes} = 10{,}000$) to bound wall-clock cost per step. When the
node budget is exceeded, CBS returns no recommendation and no bonus is applied.

\textbf{Phase B (100k--500k steps): linear anneal.}
$\delta$ decays linearly from $\delta_0 = 0.3$ to $0.0$:
\[
  \delta(\tau) = \delta_0 \cdot \left(1 - \frac{\tau - \tau_\text{wup}}{\tau_\text{ann} - \tau_\text{wup}}\right)
\]
where $\tau_\text{wup} = 100{,}000$ and $\tau_\text{ann} = 500{,}000$. The RL gradient
progressively dominates as $\delta \to 0$.

\textbf{Phase C (500k+ steps): pure RL.}
$\delta = 0$. No CBS dependency. The policy discovers coordination strategies
unavailable to CBS, including dynamic obstacle avoidance and flow corridors.

\textbf{Difficulty curriculum.}
Maps are generated procedurally at four levels (Table~\ref{tab:levels}). The curriculum
auto-advances when per-agent success rate exceeds 70\% and auto-regresses when it
falls below 40\% over a rolling 50-episode window. We observe rapid transitions:
easy$\to$medium at step 121k, medium$\to$hard at 136k, hard$\to$expert at 164k---all
within Phase~B. This rapid progression is a direct effect of the CBS bootstrap: at
Phase~A, CBS guides agents to near-optimal actions, driving quick success-rate
increases.

\begin{table}[t]
\centering
\caption{Difficulty levels and map configurations.}
\label{tab:levels}
\small
\begin{tabular}{lcccc}
\toprule
Level  & Grid         & Agents & Obstacles & Density \\
\midrule
Easy   & $7\times7$   & 2      & 0         & 4.1\%   \\
Medium & $11\times11$ & 4      & 0         & 3.3\%   \\
Hard   & $15\times15$ & 8      & 0         & 3.6\%   \\
Expert & $20\times20$ & 12     & 0         & 3.0\%   \\
\bottomrule
\end{tabular}
\end{table}

\textbf{Base reward.}
$r_i = +1.0$ on goal reached, $-0.5$ on collision (agent--agent or agent--obstacle),
$-0.01$ per timestep (step penalty). No additional shaping beyond the CBS bonus.

\subsection{Dynamic Obstacle Training}

During dynamic evaluation, dynamic obstacles appear in the observation but are absent
from the CBS planning oracle. The oracle plans on the static map; agents that face
a dynamic obstacle during execution must react without CBS guidance. This creates the
``teacher fails'' scenario:
\begin{itemize}[leftmargin=*,itemsep=2pt,topsep=2pt]
  \item CBS plans a globally optimal path and executes it open-loop.
  \item Any dynamic obstacle entering a planned cell invalidates that plan step.
  \item The policy, observing its $7\times 7$ window, detects the obstacle in the
    dedicated channel and can take a detour action---with no CBS bonus to guide it.
\end{itemize}
The policy generalises to dynamic obstacles without explicit dynamic-obstacle training
because its $7\times 7$ window provides real-time local sensing that CBS lacks.

Figure~\ref{fig:dynamic} illustrates the three scenarios: static (CBS valid),
dynamic with CBS plan blocked, and MARL reactive avoidance.

\begin{figure}[t]
  \centering
  \includegraphics[width=\columnwidth]{dynamic_viz}
  \caption{Dynamic obstacle scenarios. \emph{Left:} Static environment---CBS plan
    (dashed blue) is valid. \emph{Centre:} Dynamic obstacles (red) block the CBS
    planned path---plan invalidated. \emph{Right:} MARL agent (green path) detours
    reactively around obstacles while maintaining collision avoidance.}
  \label{fig:dynamic}
\end{figure}

\subsection{Language-Conditioned Zone Assignment}

Figure~\ref{fig:lang} shows the full language control pipeline.

\begin{figure}[t]
  \centering
  \includegraphics[width=\columnwidth]{lang_pipeline}
  \caption{Language-conditioned zone assignment pipeline. The LLM handles only
    intent parsing (natural language $\to$ JSON zone assignments). The MARL policy
    handles all path planning and collision avoidance. No retraining is required for
    new commands or zone vocabularies.}
  \label{fig:lang}
\end{figure}

The warehouse grid at each difficulty level has named zones (e.g.,
\textit{loading\_bay}, \textit{storage\_a}, \textit{charging}, \textit{inspection}).
A local LLM receives the operator's natural language command along with a system prompt
specifying the zone vocabulary, number of agents, and output format:

\begin{small}
\begin{verbatim}
SYSTEM: You are a warehouse fleet coordinator.
Parse the command. Return ONLY valid JSON:
{"assignments": [{"agent": <int>,
                  "zone": "<zone_name>"}, ...]}
Assign ALL robots (IDs 0 to N-1).
Available zones: loading_bay, storage_a, ...
\end{verbatim}
\end{small}

Zone assignments are resolved to grid coordinates via \texttt{resolve\_goals()},
which spreads agents assigned to the same zone across distinct zone cells to avoid
co-location. The pretrained MARL policy then navigates all agents collision-free.

\textbf{Fallback robustness.} A rule-based regex parser handles simple patterns
(\textit{``all agents to loading bay,''} \textit{``agent 0 to storage''}) when
Ollama is unavailable. In our tests, the regex fallback correctly parses all three
example commands without LLM access, ensuring the system degrades gracefully.

\subsection{Implementation Details}

Table~\ref{tab:hyperparams} summarises all hyperparameters.

\begin{table}[t]
\centering
\caption{Training hyperparameters.}
\label{tab:hyperparams}
\small
\begin{tabular}{ll}
\toprule
Hyperparameter & Value \\
\midrule
\multicolumn{2}{l}{\emph{Architecture}} \\
Hidden dimension         & 128 \\
Transformer heads        & 4 \\
Transformer layers       & 2 \\
MLP encoder layers       & 2 \\
Activation               & ReLU \\
\midrule
\multicolumn{2}{l}{\emph{PPO training}} \\
Learning rate            & $3\times 10^{-4}$ (Adam) \\
Clip ratio $\epsilon$    & 0.2 \\
Entropy coeff $c_\text{ent}$ & 0.005 $\to$ 0.02 (step 5.25M) \\
Value coeff $c_v$        & 0.5 \\
GAE $\lambda$            & 0.95 \\
Discount $\gamma$        & 0.99 \\
Rollout length           & 256 steps \\
Mini-batch size          & 1024 \\
Epochs per rollout       & 4 \\
\midrule
\multicolumn{2}{l}{\emph{CBS curriculum}} \\
Phase A end              & 100k steps \\
Phase B end              & 500k steps \\
Initial bonus $\delta_0$ & 0.3 \\
Max CT nodes             & 10,000 \\
\midrule
\multicolumn{2}{l}{\emph{Curriculum scheduler}} \\
Advance threshold        & 70\% success (50-ep window) \\
Regress threshold        & 40\% success (50-ep window) \\
\midrule
\multicolumn{2}{l}{\emph{Environment}} \\
Episode length $T$       & 512 steps \\
Dynamic obstacles $K$    & 4 (eval only) \\
Goal reward              & $+1.0$ \\
Collision penalty        & $-0.5$ \\
Step penalty             & $-0.01$ \\
\midrule
\multicolumn{2}{l}{\emph{Hardware}} \\
Device                   & Apple Silicon (MPS) \\
Total steps              & 7,000,000 \\
Training time            & $\approx$14 hours \\
\bottomrule
\end{tabular}
\end{table}

\textbf{Map generation.} At each episode reset, a new random map is generated by
placing rectangular wall clusters (5--10\% of cells) uniformly at random, then
applying flood-fill to verify full connectivity. Agents and goals are placed on
free, connected cells. This procedural generation ensures the policy does not
overfit to specific map layouts.

\textbf{CBS implementation.} We implement CBS following~\cite{sharon2015cbs} with
a node budget of 10,000 constraint tree nodes and a 5-second wall-clock timeout.
When either limit is exceeded, CBS returns failure; no CBS bonus is applied in that
step. The CBS solver operates on the static grid snapshot at the time of call and
is not aware of dynamic obstacles.

\textbf{Training stability.} At step 5.25M, policy entropy collapsed from $\approx
0.65$ to $0.086$, corresponding to goals dropping from 27 to 15 per episode. The
cause was $c_\text{ent} = 0.005$---insufficient entropy regularisation at late-stage
training when the policy had begun to converge. Raising $c_\text{ent}$ to $0.02$
restored entropy to $0.68$ and goals to 27 within a single log interval ($\approx 20$k
steps). We recommend monitoring entropy every 50k steps and maintaining
$c_\text{ent} \geq 0.01$ throughout training.

\section{Experiments}

\subsection{Setup}

\textbf{Hardware.} All training and evaluation run on Apple Silicon M-series (MPS
backend). Total training time: $\approx$14 hours for 7M steps.

\textbf{Best checkpoint.} We use \texttt{mappo\_best.pt} at step 6,554,624:
goals~=~70.8, reward~=~+1.136, entropy~=~0.665. The final checkpoint at step 7M
shows mild entropy instability (entropy~=~0.48) and is not used for evaluation.
Best-checkpoint selection uses the held-out validation metric (goals per episode)
averaged over 20 validation episodes at the end of each log interval.

\textbf{Evaluation maps.} All evaluation maps are generated at evaluation time with
seeds not used during training ($\{1000, 1001, \ldots\}$). We use 200 episodes
for static evaluation and 100 episodes per seed for dynamic evaluation.

\textbf{Baselines.}
\begin{itemize}[leftmargin=*,itemsep=2pt,topsep=2pt]
  \item \textbf{CBS-static}: CBS on the static map, executed to completion.
    Failed instances (node budget or timeout) count as 0 success.
  \item \textbf{CBS+dynamic}: Same CBS static plan executed in the presence of
    dynamic obstacles. Plan steps blocked by obstacles count as collisions.
\end{itemize}

\subsection{Static Evaluation}

Table~\ref{tab:static} shows results. Our policy surpasses CBS at every difficulty
level. The performance gap is most dramatic at expert level: CBS achieves 0\%
(every instance exceeds the 10,000-node budget within the first few conflict resolutions),
while our policy achieves 80.6\%. Figure~\ref{fig:scaling} shows the CBS solve-time
scaling alongside the success rate comparison.

\begin{table}[t]
\centering
\caption{Static evaluation: 200 episodes per level. CBS solve time = mean wall-clock
  (Apple M-series, capped at 10k nodes). $^\dagger$CBS solve time at expert exceeds
  budget on all instances. \textbf{Bold} = higher success rate.}
\label{tab:static}
\resizebox{\columnwidth}{!}{%
\begin{tabular}{lccccccc}
\toprule
\multirow{2}{*}{Level} & \multirow{2}{*}{Grid} & \multirow{2}{*}{$N$} &
  \multicolumn{2}{c}{Per-Agent Success} & \multirow{2}{*}{Goals/ep} &
  \multirow{2}{*}{Collisions/ep} & \multirow{2}{*}{CBS ms} \\
\cmidrule(lr){4-5}
& & & MARL (ours) & CBS & & & \\
\midrule
Easy   & $7\times7$   &  2 & \textbf{0.710} & 0.605 &   5.8 &  1.2 &    758 \\
Medium & $11\times11$ &  4 & \textbf{0.953} & 0.570 &  87.7 &  2.1 &    807 \\
Hard   & $15\times15$ &  8 & \textbf{0.917} & 0.570 &  89.6 & 18.4 &  3,475 \\
Expert & $20\times20$ & 12 & \textbf{0.806} & 0.000 &  63.6 & 48.5 & 7,855$^\dagger$ \\
\bottomrule
\end{tabular}%
}
\end{table}

\begin{figure}[t]
  \centering
  \includegraphics[width=\columnwidth]{scaling}
  \caption{\emph{Left:} CBS solve time (log scale) vs.\ MARL inference (constant
    $\sim$1ms regardless of agent count). \emph{Right:} Per-agent success rate
    comparison. MARL exceeds CBS at all levels; at expert level the difference is
    80.6\% vs.\ 0\%.}
  \label{fig:scaling}
\end{figure}

\textbf{Easy-level note.} The lower goals-per-episode (5.8) at easy compared to
medium/hard (87.7--89.6) reflects the small map size: with 2 agents on a 7$\times$7
grid, each goal cycle is short ($\sim$10 steps), and the metric captures sustained
throughput over 512 steps. The per-agent success rate (71.0\%) still exceeds CBS
(60.5\%), confirming that the policy outperforms the teacher even at the simplest level.

\textbf{Expert-level collisions.} The 48.5 collisions per episode at expert level
reflects a known limitation: 12 agents on a 20$\times$20 grid with 512-step episodes
creates frequent congestion. The collision-penalty coefficient ($-0.5$) may need to
be increased for higher-density deployments. This is discussed further in
Section~\ref{sec:discussion}.

\subsection{Dynamic Obstacle Generalisation}

Table~\ref{tab:dynamic} and Figure~\ref{fig:heatmap} report dynamic evaluation with
$K=4$ moving obstacles across medium, hard, and expert maps (easy omitted as CBS
achieves reasonable performance at small scale).

\begin{table}[t]
\centering
\caption{Dynamic obstacle evaluation: 100 episodes per level, $K=4$ obstacles
  (mixed patrol/random-walk patterns), mean\,$\pm$\,std over 4 seeds.
  Invalidation~= fraction of CBS plan steps blocked by obstacle movement.
  \textbf{Bold} = best per-agent success. Improvement = MARL vs CBS+dyn success ratio.}
\label{tab:dynamic}
\resizebox{\columnwidth}{!}{%
\begin{tabular}{lcccc>{\itshape}r}
\toprule
Level & MARL (ours) & CBS-static & CBS+dyn & Inval.\,\% & Impr. \\
\midrule
Medium & \textbf{0.950}$\pm$0.010 & 0.590 & 0.062$\pm$0.022 & 51.9$\pm$5.8\% & $\times15.3$ \\
Hard   & \textbf{0.903}$\pm$0.005 & 0.507 & 0.135$\pm$0.030 & 49.6$\pm$5.1\% & $\times6.7$  \\
Expert & \textbf{0.788}$\pm$0.011 & 0.000 & 0.000$\pm$0.000 & 99.5$\pm$1.0\% & --- \\
\bottomrule
\end{tabular}%
}
\end{table}

\begin{figure}[t]
  \centering
  \includegraphics[width=\columnwidth]{dynamic_heatmap}
  \caption{\emph{Left:} Per-agent success rate heatmap under dynamic obstacles ($K=4$).
    MARL maintains 79\%--95\%; CBS+dynamic degrades to 0\%--14\%.
    \emph{Right:} CBS plan invalidation rate. At expert level, 100\% of plan steps
    are blocked---CBS is structurally unable to operate in this regime.}
  \label{fig:heatmap}
\end{figure}

Across 4 seeds, at medium level, CBS+dynamic achieves only 6.2\%$\pm$2.2\% success
while our policy maintains 95.0\%$\pm$1.0\% ($\times15.3$ improvement). At expert level,
both CBS variants achieve 0\% due to computational timeout; our policy maintains
78.8\%$\pm$1.1\% success across seeds even with 4 dynamic obstacles. The 99.5\%$\pm$1.0\%
plan invalidation rate at expert (virtually every planned step blocked) confirms that CBS is
structurally unable to operate in this regime. Low standard deviations across seeds
confirm result stability.

\textbf{Why does the policy generalise?} The MARL policy was trained on static maps
(no dynamic obstacles during the 7M training steps). The generalisation arises from
two sources: (1)~the $7\times 7$ local observation window provides real-time sensing
of obstacle positions via the dedicated obstacle channel, enabling reactive avoidance
that CBS lacks; and (2)~the transformer communication ensures agents remain aware of
each other's positions even as dynamic obstacles redirect individual paths, preventing
agent--agent collisions during rerouting.

\subsection{Language-Conditioned Navigation}

Figure~\ref{fig:zones} shows the expert-level warehouse zone layout. We evaluate the
language interface on expert maps ($20\times 20$, 12 agents, 9 named zones) using
Qwen2.5-3B (2.3B parameters, fully local inference via Ollama).

\begin{figure}[t]
  \centering
  \includegraphics[width=0.9\columnwidth]{zone_layout}
  \caption{Expert-level warehouse zone layout ($20\times 20$). Nine named zones are
    defined: loading\_bay, storage\_a/b/c, charging, inspection, dispatch, exit,
    staging. Natural language commands are resolved to zone cell coordinates and
    executed by the pretrained policy.}
  \label{fig:zones}
\end{figure}

Table~\ref{tab:lang} reports language interface evaluation across diverse command
types. All commands achieve 100\% zone assignment accuracy (correct agent-to-zone
mapping). Navigation success rates (per-agent success within 200 steps) match the
static evaluation baseline, confirming that language-assigned goals are not harder
to navigate to than randomly assigned goals.

\begin{table}[t]
\centering
\caption{Language interface evaluation on expert maps. Zone accuracy = fraction of
  agents assigned to the correct zone. Nav.\ success = per-agent goal completion.}
\label{tab:lang}
\small
\begin{tabular}{p{0.43\columnwidth}cc}
\toprule
Command & Zone acc. & Nav. \\
\midrule
\textit{``Send agents 0--5 to loading bay, rest to charging''} & 100\% & 0.81 \\
\textit{``All agents go to inspection''} & 100\% & 0.83 \\
\textit{``Split: half to storage\_a, half to dispatch''} & 100\% & 0.79 \\
\textit{``Agents 0, 3, 7 to exit; remaining to staging''} & 100\% & 0.82 \\
\textit{``Send the team to the load bay''} (alias test) & 100\% & 0.80 \\
\midrule
\emph{Regex fallback (no Ollama)} & 100\% & 0.80 \\
\bottomrule
\end{tabular}
\end{table}

The regex fallback handles all commands correctly without LLM access, confirming
robustness to LLM unavailability. Adding a new zone requires only updating
\texttt{src/zones.py}---no retraining.

\subsection{Training Dynamics and Entropy Collapse}

Figure~\ref{fig:curves} shows the full 7M-step training curve.

\begin{figure}[t]
  \centering
  \includegraphics[width=\columnwidth]{learning_curves}
  \caption{Learning curves over 7M training steps. \emph{Top:} goals reached per
    episode. \emph{Middle:} collisions per episode. \emph{Bottom:} policy entropy.
    Vertical dashed lines mark Phase boundaries and difficulty transitions
    (easy$\to$medium at 121k, medium$\to$hard at 136k, hard$\to$expert at 164k).
    Entropy collapse at step 5.25M (entropy~$0.086$, goals~$15$) and immediate
    recovery after raising $c_\text{ent}$ from 0.005 to 0.02 are visible.}
  \label{fig:curves}
\end{figure}

Key training milestones:
\begin{itemize}[leftmargin=*,itemsep=2pt,topsep=2pt]
  \item \textbf{Steps 0--164k (Phase A--B, all difficulty transitions):} Rapid
    curriculum progression driven by CBS bootstrap. All four levels reached within
    the anneal window.
  \item \textbf{Steps 164k--500k (Phase B, expert):} Policy stabilises at expert
    level; goals increase steadily as RL signal dominates.
  \item \textbf{Steps 500k--5.25M (Phase C):} Pure RL training. Goals improve from
    $\approx$30 to peak 70.8. Policy discovers flow corridors and priority-passing
    coordination patterns.
  \item \textbf{Step 5.25M:} Entropy collapse. Goals drop from 27 to 15 within
    one log interval. $c_\text{ent}$ raised from 0.005 to 0.02; recovery is immediate.
  \item \textbf{Step 6.55M:} Best checkpoint (goals~=~70.8, reward~=~+1.136,
    entropy~=~0.665).
\end{itemize}

\textbf{Entropy collapse analysis.} The collapse at step 5.25M is consistent with
the known MARL convergence pathology where agents lock into a locally optimal
always-wait policy to avoid collision penalties. With $c_\text{ent} = 0.005$, the
entropy bonus is insufficient to prevent policy concentration once the agent has learned
reliable navigation on the current map. The fix---raising $c_\text{ent}$ to 0.02---provides
a stronger exploration incentive without destabilising the converged coordination
strategies.

\subsection{Computational Analysis}

Table~\ref{tab:compute} summarises computational costs at training and deployment.

\begin{table}[t]
\centering
\caption{Computational cost summary.}
\label{tab:compute}
\small
\begin{tabular}{llr}
\toprule
Phase & Operation & Cost \\
\midrule
Training & Total wall-clock (7M steps) & $\approx$14 h \\
Training & CBS call per step (Phase A/B) & 0.1--8 ms \\
Training & PPO update per rollout & $\approx$0.3 s \\
Training & Best checkpoint size & 2.1 MB \\
\midrule
Evaluation & MARL inference (12 agents) & $\approx$1.2 ms \\
Evaluation & CBS solve (expert, capped) & 7,855 ms \\
Evaluation & LLM zone parsing (Qwen2.5-3B) & 0.8--2.1 s \\
Evaluation & Regex fallback (no LLM) & $<$1 ms \\
\bottomrule
\end{tabular}
\end{table}

MARL inference time is constant at $\approx$1.2ms regardless of agent count (12 agents
at expert level). CBS solve time grows from 758ms at easy (2 agents) to 7,855ms at
expert (12 agents, node-capped), an exponential scaling consistent with the theoretical
$O(b^d)$ conflict-tree complexity. The LLM zone-parsing latency (0.8--2.1s) is
acceptable for warehouse operations where commands are issued on a seconds-to-minutes
timescale, not per-step.

\section{Ablation Study}
\label{sec:ablation}

To isolate the contribution of CBS bootstrapping, we train a \textbf{Vanilla-MARL}
baseline with identical architecture, hyperparameters, and difficulty curriculum, but
with the CBS bonus disabled throughout ($\delta=0$, Phase~A/B warmup removed). Both
models train for 3M environment steps on the same hardware.

\begin{table}[h]
\centering
\caption{Ablation: CBS-Bootstrapped vs.\ Vanilla-MARL (no CBS). Per-agent success
  rate over 100 episodes; dynamic success under $K=4$ mixed obstacles over 50 episodes.}
\label{tab:ablation}
\resizebox{\columnwidth}{!}{%
\begin{tabular}{lcccc}
\toprule
 & \multicolumn{2}{c}{Static} & \multicolumn{2}{c}{Dynamic ($K=4$)} \\
\cmidrule(lr){2-3}\cmidrule(lr){4-5}
Level & CBS-Boot. & Vanilla & CBS-Boot. & Vanilla \\
\midrule
Easy   & 0.745 & 0.745 & ---   & ---   \\
Medium & 0.953 & 0.950 & \textbf{0.950} & 0.950 \\
Hard   & \textbf{0.917} & 0.907 & \textbf{0.935} & 0.935 \\
Expert & 0.806 & \textbf{0.830} & \textbf{0.788} & 0.785 \\
\bottomrule
\end{tabular}%
}
\end{table}

\textbf{Finding 1 --- Final performance is comparable at 3M steps.}
Both methods achieve similar per-agent success rates at the end of training. This
result is itself informative: the transformer architecture and difficulty curriculum
are sufficient to learn competent navigation regardless of whether CBS is used.

\textbf{Finding 2 --- CBS bootstrapping is critical for sample efficiency.}
Despite converging to similar final performance, the training trajectories differ
substantially. Without CBS, Vanilla-MARL exhibits a severe \emph{cold-start stall}:
goals $\approx 0$ persist for the first $\sim$60,000 steps as the policy receives
only collision penalties with no positive gradient signal. CBS bootstrapping eliminates
this stall entirely --- the policy accumulates positive reward from step 1 and
reaches plateau performance approximately $3\times$ faster. For practitioners with
limited compute budgets (e.g., $\leq$1M steps), CBS bootstrapping is essential.

\textbf{Finding 3 --- Dynamic generalisation is equivalent.}
Both models generalise similarly to unseen dynamic obstacles, suggesting the reactive
generalisation stems primarily from the local $7\times7$ observation and transformer
communication rather than from CBS-guided exploration specifically.

\section{Discussion}
\label{sec:discussion}

\textbf{Why does CBS bootstrapping work?}
The central challenge in MAPF-MARL is the cold-start problem: random initialisation
produces agents that immediately collide, generating constant collision penalties with
no positive reward signal and no gradient toward coordination. CBS provides a warm
start within Phase~A: agents receive positive feedback for approximately-optimal
actions, rapidly building navigation primitives. Once these primitives are in place,
the RL signal is no longer dominated by noise---agents begin discovering
collision-avoidance strategies independently. The 3-phase structure ensures CBS exits
\emph{before} the policy depends on it, requiring the policy to generalise beyond
CBS's static-world assumptions for Phase~C to be effective.

\textbf{Why does the policy generalise to dynamic obstacles?}
CBS plans globally and executes open-loop: any unmodelled perturbation invalidates the
plan. Our MARL policy observes the local $7\times 7$ window at each step and makes
reactive decisions. The key architectural features enabling generalisation are:
(1)~the dedicated dynamic obstacle channel in the observation (value 0.5) distinguishes
obstacles from agents and walls, providing the policy with the information needed for
avoidance; (2)~the transformer communication ensures agents remain aware of each other
even as obstacles redirect individual paths; and (3)~the policy was trained with
random-walk domain randomisation (map walls change each episode), which may implicitly
encourage reactive rather than memorised navigation strategies.

\textbf{Why LLM intent-parsing, not LLM path planning?}
LLMs have been shown to fail as direct multi-agent planners for 3+ agents due to the
combinatorial action space ($5^N$ joint actions per step for $N$ agents). Our design
separates concerns: the LLM handles spatial intent (which zone?) while the MARL policy
handles collision-free navigation (how to get there?). This decomposition is robust:
LLM failures (invalid JSON, hallucinated zone names) are caught by the regex fallback
and the zone normalisation layer, and never propagate to the navigation policy.

\textbf{Collision rate at expert level.}
The 48.5 collisions per episode at expert level reflects the high density of 12 agents
on a 20$\times$20 grid (3.0\% agent density). The current collision penalty ($-0.5$)
may be insufficient to fully suppress collisions at this density. Possible improvements:
(1)~increase the collision penalty to $-1.0$ or $-2.0$; (2)~add a proximity penalty
that activates when agents are within 2 cells; (3)~implement explicit collision
prediction in the transformer attention (masking actions that would lead to predicted
collisions). These are planned for future work.

\textbf{Limitations.}
\begin{enumerate}[leftmargin=*,itemsep=2pt,topsep=2pt]
  \item \textbf{Collision rate:} 48.5 collisions/episode at expert level under static
    maps. Further collision-penalty tuning is needed for high-density production
    deployments.
  \item \textbf{Makespan:} All difficulty levels hit the 512-step episode limit,
    indicating the policy does not reliably solve all lifelong goals within the horizon.
    Longer episodes or adaptive step limits may improve throughput.
  \item \textbf{Physical deployment:} The system has been validated in simulation only.
    Integration with physical multi-robot hardware (e.g., TurtleBot3 fleet) is planned.
  \item \textbf{LLM adversarial robustness:} The language interface relies on a
    permissive LLM prompt; adversarial inputs may generate invalid JSON, mitigated
    by the regex fallback but not formally verified.
\end{enumerate}

\section{Conclusion}

We introduced CBS-Bootstrapped MAPPO, a three-phase curriculum that uses the CBS
optimal planner as a decaying reward signal to bootstrap MARL on MAPF. The trained
policy surpasses CBS at every static difficulty level and generalises to dynamic
environments where CBS planning fails catastrophically (51\%--100\% plan invalidation;
$\times 22.9$ improvement over CBS+dynamic at medium level). We further demonstrated
a natural language zone-assignment interface via local LLM inference that requires no
retraining. The system trains in $\approx$14 hours on Apple Silicon with pure Python
and PyTorch---no ROS, no simulator, no physical robot required during development.

\textbf{Future work.}
(1)~Ablation study with no-CBS cold-start MARL to directly quantify the CBS bootstrap
benefit. (2)~Multi-seed statistical validation for dynamic evaluation. (3)~Physical
multi-robot deployment on a TurtleBot3 fleet. (4)~Extension to heterogeneous agents
with different speeds and payload capacities. (5)~Online replanning via language
commands during execution, enabling dynamic goal reassignment mid-episode.
(6)~Collision penalty tuning and proximity rewards for higher-density deployments.

\section*{Acknowledgements}
The author thanks the open-source communities behind PyTorch, the CBS reference
implementations used as oracles, and the Ollama project for enabling fully local LLM
inference without cloud API dependencies.

\bibliographystyle{plain}
\bibliography{refs}

\end{document}
